\title{Direct Preference Optimization: Your Language Model is Secretly a Reward Model}

\begin{abstract}
Although large-scale unsupervised language models (LMs) acquire extensive general knowledge and exhibit some reasoning capabilities, precisely steering their outputs remains challenging because their training lacks supervision. To introduce such controllability, current approaches typically involve gathering human annotations indicating the relative merits of different model outputs and then further training the LM so its behavior aligns with these annotated preferences—commonly through reinforcement learning from human feedback (RLHF). However, RLHF is often intricate and unstable, involving an initial phase where a reward model is trained to encapsulate human preferences, followed by reinforcement learning to adapt the LM to maximize the learned reward while avoiding significant deviation from the original parameters. In this work, we propose an alternative parameterization of the reward model used in RLHF, which makes it possible to analytically derive the optimal policy in closed form. This allows us to resolve the typical RLHF optimization problem simply with a classification loss. We refer to this approach as \textit{Direct Preference Optimization} (DPO). DPO provides a robust, efficient, and low-overhead alternative, obviating the need for language model sampling during fine-tuning and requiring little hyperparameter tuning. Through empirical evaluation, we demonstrate that DPO is capable of fine-tuning LMs for preference alignment at least on par with, and sometimes superior to, established techniques. Importantly, DPO surpasses RLHF methods based on PPO in managing generative sentiment, and attains equivalent or improved results in tasks like summarization and single-turn dialogue—while being considerably more straightforward to deploy and train.
\end{abstract}

\section{Introduction}
Large-scale unsupervised language models (LMs) trained on extensive corpora often exhibit remarkable emergent properties~\citep{chowdhery2022palm, brown2020language, touvron2023llama, bubeck2023sparks}. Yet, these models learn from human-generated data that encompasses a vast spectrum of intentions, competencies, and objectives. Not all of these are suitable for emulation; for instance, while it is beneficial for an AI coding assistant to \textit{recognize} frequent programming errors so as to remedy them, it is preferable for the system’s code generation to reflect the exceptional—albeit uncommon—exemplars of programming expertise that exist in the training distribution. Likewise, we desire that the language model be \textit{informed} about prevalent misconceptions, even if 50\% of individuals hold them, but we unequivocally do not wish the model to endorse such misconceptions as valid in half of the relevant interactions. Put differently, isolating the \emph{preferred outputs and conduct} from the broader spectrum of the model’s \textit{knowledge and competencies} is essential for creating AI that is robust, safe, and manageable \citep{ouyang2022training}. Whereas predominant techniques align language model outputs with human preferences using reinforcement learning (RL), we demonstrate that the standard RL-derived objective underpinning such techniques can, in fact, be realized directly and precisely by applying a basic binary cross-entropy loss—significantly streamlining the process of preference alignment.

In broad terms, prevalent approaches introduce targeted behavioral dispositions into language models by leveraging carefully constructed datasets that capture human-preferred actions and attitudes deemed helpful or safe. This phase of preference adaptation follows an initial unsupervised pre-training on expansive text datasets. While the most naive method for this adaptation consists of supervised fine-tuning using examples of high-quality human-authored outputs, the preeminent paradigm is reinforcement learning from human (or AI-generated) feedback, known as RLHF/RLAIF~\citep{christiano2017deep,bai2022constitutional}. RLHF methodologies involve training a reward model based on human preference data, and subsequently employing RL to optimize the model policy so that it produces responses ranked highly by the reward model, all while maintaining close adherence to the pre-trained model’s behavior. Although RLHF-based models yield notable advances in dialogue and programming tasks, their learning pipeline is far more intricate compared to supervised techniques: it requires training several separate models, policy sampling within a training loop, and incurs notable computational overhead.

This work demonstrates a method for optimizing language models to better reflect human preferences, eliminating the need for explicit reward models or reinforcement learning approaches. We introduce \textit{Direct Preference Optimization} (DPO), a novel technique that implicitly targets the same reward maximization objective under a KL-divergence constraint as conventional RLHF algorithms, yet is more straightforward both in terms of implementation and training. Conceptually, DPO functions by boosting the ratio of log probabilities between preferred and non-preferred outputs, while employing a flexible, per-instance importance weighting scheme to counteract the degradation observed with simpler probability ratio objectives. Similar to prior methods, DPO utilizes a theoretical model of preference (e.g., the Bradley-Terry model; \cite{bradley1952rankanalysis}) to assess the correspondence between reward functions and actual human preference data. However, in contrast to prior approaches that leverage the preference model first to train a reward function, and subsequently to optimize a policy against this learned reward, DPO performs a reparameterization, enabling direct formulation of the preference loss in terms of the policy itself. This reparameterization allows DPO to employ a standard binary cross entropy loss on datasets comprising pairwise human preference judgments, facilitating direct policy optimization and yielding an optimal policy with respect to an implicit reward derived from preference annotations.

The primary contribution of this study is the presentation of Direct Preference Optimization (DPO), an intuitive, reinforcement-learning-free framework for learning language models from human preference data. Empirical results indicate that DPO matches or outperforms established techniques—including RLHF using PPO—across a variety of tasks such as sentiment control, summarization, and conversation, with models scaling up to 6B parameters.

\section{Related Work}

Progressively larger self-supervised language models have demonstrated the ability to handle certain tasks without any fine-tuning (zero-shot) \citep{radford2019language}, or by leveraging just a handful of examples provided in prompts (few-shot) \citep{gpt3,megatron,chowdhery2022palm}. Nonetheless, their efficacy on specific downstream tasks and alignment with users’ expectations can be markedly enhanced when these models undergo additional fine-tuning on corpora containing explicit instructions and human-crafted completions \citep{mishra-etal-2022-cross,sanh2022multitask,chung2022scaling,thoppilan2022lamda}. This strategy, termed `instruction-tuning,’ equips LLMs with the capacity to extrapolate to novel instructions not encountered during fine-tuning, thereby broadening their practical utility \citep{chung2022scaling}. In practice, gathering human preference comparisons of model responses is often less demanding than curating expert-level demonstrations; therefore, recent approaches have prioritized refining LLMs using such preference-based datasets. This shift has fostered progress in translation \citep{kreutzer-etal-2018-reliability}, summarization \citep{stiennon2022learning,ziegler2020finetuning}, narrative generation \citep{ziegler2020finetuning}, and following user instructions \citep{ouyang2022training,ramamurthy2023is}. Typically, these methods involve training a neural reward model to reflect user preferences, frequently adopting a framework such as the Bradley-Terry model \citep{bradley1952rankanalysis}, and then applying reinforcement learning—most commonly REINFORCE \citep{williams1992reinforce}, proximal policy optimization (PPO; \cite{schulman2017proximal}), or their derivatives \citep{ramamurthy2023is}—to adapt the language model toward maximizing the reward. Parallel research directions use instruction-tuned LLMs informed by human feedback to autonomously generate new preference data, emphasizing particular qualities like safety or non-maliciousness \citep{bai2022constitutional}. In these setups, human supervision may be limited to providing textual guidelines that guide model annotation. Such developments represent a convergence between two major research areas: one focused on employing reinforcement learning to train language models for diverse goals~\citep{Ranzato2015SequenceLT,paulus2018a,wu2018learning}, and the other on generic algorithms for learning from human preferences \citep{christiano2017deep,kupcsik2018learning}. Although harnessing relative human judgments is attractive, implementing reinforcement learning for large-scale language model fine-tuning introduces considerable complexity; the present study offers a principled alternative for preference optimization that circumvents reliance on RL.

Beyond linguistic applications, the study of learning policies from preferences has been explored within both bandit frameworks and reinforcement learning, leading to various methodologies. In contextual bandits, learning from preferences or ordinal rankings instead of numerical rewards is termed the contextual dueling bandit (CDB) problem \citep{yue2012karmed,dudik2015contextual}. Since CDBs operate without absolute reward signals, theoretical work replaces the notion of an optimal policy with that of a \textit{von Neumann winner}—a policy expected to outperform any competitor at least half the time \citep{dudik2015contextual}. Notably, CDBs presume the availability of preference information in an online fashion, contrasting with learning from human preferences, where typically a static dataset of annotated preference pairs is employed \citep{yan2022human}. Likewise, \textit{preference-based RL} (PbRL) eschews rewards in favor of binary preference signals inferred from an unknown "scoring" function \citep{BusaFekete2014,ruiz2023dueling}. Numerous PbRL algorithms have been developed, some allowing the use of offline preferences, though most first estimate this latent scoring (or reward) function before optimizing policies with respect to it \citep{jain2013learning,BusaFekete2014,christiano2017deep,sadigh2017active,kupcsik2018learning}. By contrast, our contribution is to introduce a direct policy optimization strategy that seeks to meet preference criteria in a single stage.

\section{Preliminaries}\label{section:prelims}

We briefly outline the RLHF procedure introduced in \citeauthor{ziegler2020finetuning} and followed in subsequent works \citep{stiennon2022learning, bai2022training, ouyang2022training}, typically composed of three stages: 1) supervised fine-tuning (SFT); 2) preference data collection and reward model training; and 3) reinforcement learning-based optimization.

\textbf{SFT}: RLHF pipelines are initiated by further training a large pre-trained language model using supervised learning on curated, task-specific datasets (such as dialogue or summarization data) to produce a supervised model, denoted $\pi^\text{SFT}$.

\textbf{Reward Modelling Phase}: Next, the SFT model is prompted on inputs $x$ to generate candidate answer pairs $(y_1, y_2)\sim \pi^\text{SFT}(y \mid x)$. Human annotators are then asked to indicate their preference between the two, i.e., $y_w\succ y_l \mid x$, where $y_w$ and $y_l$ indicate the favored and less favored completions, respectively, from the sampled pair. These preferences are assumed to reflect an unobserved reward function $r^*(y, x)$ to which we lack direct access. Several statistical models exist for capturing preference information, with the Bradley-Terry (BT) \cite{bradley1952rankanalysis} model being a standard choice (though the broader Plackett-Luce ranking framework \citep{plackett1975analysis, luce2012individual} is also suitable, particularly when multiple rankings are available). Under the BT model, the distribution of human preferences $p^*$ is modeled as:

\begin{equation}\label{eq:bradley-terry}
    p^*(y_1\succ y_2 \mid x)=\frac{\exp\left(r^*(x, y_1)\right)}{\exp\left(r^*(x, y_1)\right) + \exp\left(r^*(x, y_2)\right)}.
\end{equation}

Given a fixed dataset of preference comparisons, denoted $\mathcal{D} = \bigl\{x^{(i)}, y_w^{(i)}, y_l^{(i)}\bigr\}_{i=1}^N$, drawn according to $p^*$, we may introduce a parameterized reward function $r_{\phi}(x, y)$ and learn the parameters via maximum likelihood estimation. When casting this problem as binary classification, the resulting objective becomes the negative log-likelihood loss:

\begin{equation}\label{eq:reward_model}
    \mathcal{L}_R(r_{\phi}, \mathcal{D}) = -\mathbb{E}_{(x, y_w, y_l)\sim \mathcal{D}}\bigl[\log \sigma(r_{\phi}(x, y_w)- r_{\phi}(x, y_l))\bigr]
\end{equation}

where $\sigma$ denotes the logistic sigmoid. In language modeling applications, $r_{\phi}(x, y)$ is generally initialized from the SFT policy $\pi^\text{SFT}(y \mid x)$ and augmented with a linear projection applied to the transformer’s final hidden states to yield a scalar reward prediction \cite{ziegler2020finetuning}. To reduce reward variance, it is common practice to normalize the outputs so that $\mathbb{E}_{x,y\sim \mathcal{D}}\left[r_\phi(x, y)\right] = 0$ for every $x$.

\textbf{RL Fine-Tuning Phase}: Once the reward model is trained, it is employed as the feedback mechanism for further language model fine-tuning via RL. Following the conventions of previous research~\citep{jaques2017sequence, jaques2020human}, the corresponding optimization problem is given by

\begin{equation}\label{eq:RL}
\max_{\pi_{\theta}}  \mathbb{E}_{x\sim \mathcal{D}, y\sim \pi_{\theta}(y \mid x)}\bigl[r_{\phi}(x, y)\bigr] - \beta\mathbb{D}_{\textrm{KL}}\bigl[\pi_{\theta}(y\mid x)\mid \mid \pi_\text{ref}(y\mid x)\bigr],
\end{equation}

with $\beta$ as a regularization coefficient controlling divergence from the reference policy $\pi_\text{ref}$—in this context, usually the initial SFT model $\pi^\text{SFT}$. 
Practically, $\pi_\theta$ is also initialized from $\pi^\text{SFT}$. The KL term is crucial, as it constrains the policy to remain close to regions where the reward model’s predictions remain valid and guards against generation collapse, thus preserving answer diversity. Because text generation is inherently discrete and non-differentiable, standard gradient-based methods cannot be applied; thus, optimization typically proceeds with reinforcement learning. The dominant approach \citep{ziegler2020finetuning, stiennon2022learning, bai2022training, ouyang2022training} reformulates the scalar reward as ${r(x, y) = r_{\phi}(x, y) -\beta (\log \pi_{\theta}(y\mid x) - \log \pi_\text{ref}(y\mid x))}$ and employs proximal policy optimization (PPO) \cite{schulman2017proximal} to optimize it. 

\section{Direct Preference Optimization}\label{sec:DPO}

Given the obstacles that arise when employing reinforcement learning algorithms for large-scale language model fine-tuning, we seek a principled, streamlined method to directly optimize policies from preference data. In contrast to conventional RLHF pipelines—which independently train a reward function, then use reinforcement learning to maximize it—our method utilizes a specific reward parameterization, making it possible to derive the optimal policy in closed form, bypassing the need for an RL loop. We will elaborate on this point below. The crucial insight underlying our approach is that the mathematical relationship between reward models and their associated optimal policies enables the transformation of losses formulated over rewards into policy-based objectives. By employing this reparameterization, we forgo fitting an explicit reward model while still adhering to established models of human preferences, such as the Bradley-Terry formulation.

\textbf{Reformulation of the DPO objective.} Our derivation begins from the conventional RL objective described in Eq.~\ref{eq:RL}, parameterized by a general reward function $r$. Previous studies~\citep{peters2007reinforcement, peng2019advantage, korbak2022reinforcement, go2023aligning} have established that the solution optimizing the KL-regularized expected reward as specified in Eq.~\ref{eq:RL} adopts the structure:

\begin{equation}\label{eq:op_policy}
    \pi_r(y\mid x) = \frac{1}{Z(x)}\pi_\text{ref}(y\mid x)\exp\left(\frac{1}{\beta}r(x, y)\right),
\end{equation}

with $Z(x) =\sum_{y}\pi_\text{ref}(y\mid x)\exp\left(\frac{1}{\beta}r(x, y)\right)$ serving as a normalization constant (partition function). A detailed derivation is provided in Appendix \ref{app:derivation1}. While it is possible to substitute the MLE estimator $r_{\phi}$ in place of the true reward function $r^*$, direct computation of $Z(x)$ remains computationally burdensome~\citep{korbak2022reinforcement, go2023aligning}, limiting practical application of this formulation. Nevertheless, it is possible to manipulate Eq.~\ref{eq:op_policy} to resolve the reward as a function of the induced optimal policy $\pi_r$, the baseline policy $\pi_\text{ref}$, and $Z(\cdot)$. To accomplish this, one can logarithmically transform Eq.~\ref{eq:op_policy} and rearrange algebraically, yielding:

\begin{equation}\label{eq:main_eq}
    r(x,y) =\beta \log \frac{\pi_r(y\mid x)}{\pi_\text{ref}(y\mid x)} + \beta \log Z(x).
\end{equation}

Applying this result to both the true reward $r^*$ and its associated optimal policy $\pi^*$, and leveraging the fact that the Bradley-Terry model considers only pairwise reward differences—specifically, ${p^*(y_1 \succ y_2 \mid x) = \sigma(r^*(x, y_1) - r^*(x, y_2))}$—one can substitute the expression for $r^*(x,y)$ from Eq.~\ref{eq:main_eq} into the preference probability formula (Eq.~\ref{eq:bradley-terry}). This substitution eliminates the partition function, so the human preference likelihood may be formulated purely in terms of the optimal policy $\pi^*$ and reference $\pi_\text{ref}$. Thus, the preference relationship under the optimal RLHF policy $\pi^*$ within the Bradley-Terry framework is as follows:

\begin{equation}\label{eq:objective}
    p^*(y_1\succ y_2 \mid x)=\frac{1}{1 + \exp\left(\beta \log \frac{\pi^*(y_2\mid x)}{\pi_\text{ref}(y_2\mid x)} - \beta \log \frac{\pi^*(y_1\mid x)}{\pi_\text{ref}(y_1\mid x)}\right)}
\end{equation}

The full derivation is available in Appendix~\ref{app:derivation2}. While this presentation utilizes the Bradley-Terry framework, parallel results can be obtained for broader Plackett-Luce models~\citep{plackett1975analysis, luce2012individual}, with derivations in Appendix~\ref{app:plackett_luce_models}.

With this representation of human preference probabilities in terms of the policy, we are now equipped to introduce a maximum likelihood approach for a parametric policy $\pi_\theta$. Drawing an analogy with reward modeling (as in Eq.~\ref{eq:reward_model}), the loss function for the policy is defined as:

\begin{equation}\label{eq:optimum_model}
    \mathcal{L}_\text{DPO}(\pi_{\theta}; \pi_\text{ref}) = -\mathbb{E}_{(x, y_w, y_l)\sim \mathcal{D}}\left[\log \sigma \left(\beta \log \frac{\pi_{\theta}(y_w\mid x)}{\pi_\text{ref}(y_w\mid x)} - \beta \log \frac{\pi_{\theta}(y_l\mid x)}{\pi_\text

Through this approach, we approximate an implicit reward function by employing a distinct parameterization, for which the corresponding optimal policy directly corresponds to $\pi_\theta$. Furthermore, since this methodology is formally equivalent to fitting a reparameterized Bradley-Terry model, it inherits desirable theoretical characteristics, such as consistency under appropriate assumptions about the preference data distribution \cite{bong2022generalized}. Further exploration of the theoretical aspects of DPO in the context of related research is presented in Section~\ref{sec:theory}.

\textbf{Mechanistic role of the DPO update.} To gain insight into how DPO operates, it is instructive to examine the gradient of its loss function, $\mathcal{L}_\text{DPO}$. The derivative with respect to $\theta$ can be expressed as:

\begin{multline*}\label{eq:gradient}
    \nabla_\theta \mathcal{L}_\text{DPO}(\pi_\theta;\pi_\text{ref}) = \\ -\beta\mathbb{E}_{(x, y_w, y_l) \sim \mathcal{D}} \bigg[\underbrace{\sigma(\hat{r}_\theta(x, y_l) - \hat{r}_\theta (x, y_w))}_\text{larger contribution if the reward prediction is incorrect}\bigg[\underbrace{\nabla_\theta\log \pi(y_w \mid x)}_\text{boost probability of $y_w$} - \underbrace{\nabla_\theta\log\pi(y_l \mid x)}_\text{reduce probability of $y_l$}\bigg]\bigg],
\end{multline*}

where the implicit reward $\hat{r}_\theta(x, y)$ is defined as $\beta \log \frac{\pi_\theta(y \mid x)}{\pi_\text{ref}(y \mid x)}$, relating the language model $\pi_\theta$ to a fixed reference $\pi_\text{ref}$ (see further discussion in Section~\ref{sec:theory}). Conceptually, this gradient incentivizes the model to increase the likelihood of preferred responses $y_w$ while lowering the probability assigned to less-preferred alternatives $y_l$. Critically, the influence of each example is weighted according to the extent to which the implicit reward function $\hat{r}_\theta$ erroneously favors the less-desirable completion, with scaling governed by $\beta$, which is also tied to the KL regularization strength. Empirical results underscore the necessity of this weighting; omitting it, as in a simple unweighted variant, can lead to language model collapse (see Appendix Table~\ref{tab:unlikelihood_generations}).

\textbf{DPO overview.} The typical DPO workflow proceeds as follows: First, for each prompt $x$, completions $y_1, y_2 \sim \pi_\text{ref}(\cdot \mid x)$ are sampled and annotated with human preference data to construct the offline preferences dataset $\mathcal{D} = \{x^{(i)}, y_w^{(i)}, y_l^{(i)}\}_{i=1}^N$. Next, the language model $\pi_\theta$ is trained by minimizing $\mathcal{L}_\text{DPO}$, conditioned on the specific choices of $\pi_\text{ref}$, $\mathcal{D}$, and $\beta$. 

In practical scenarios, preference datasets that have already been curated and are publicly accessible are preferred over generating new samples and collecting human annotations. Because such datasets are typically generated using $\pi^\text{SFT}$, it is standard to initialize $\pi_\text{ref}$ with $\pi^\text{SFT}$ if it is available. In instances where $\pi^\text{SFT}$ cannot be used, $\pi_\text{ref}$ is initialized by maximizing the likelihood of preferred completions ${(x, y_w)}$ over the dataset, i.e., ${\pi_\text{ref} = \argmax_{\pi}\mathbb{E}_{x, y_w \sim \mathcal{D}}\left[\log \pi(y_w \mid x)\right]}$. This approach helps reduce the mismatch between the intractable true reference distribution and the approximation $\pi_\text{ref}$ employed in DPO. Appendix~\ref{app:implementation} contains additional implementation details and hyperparameter configurations.

\section{Theoretical Analysis of DPO}

This section deepens the theoretical understanding of DPO, interpreting its methodology, furnishing analytical justification, and connecting its strengths with known limitations of actor-critic methods, such as those found in RLHF paradigms employing PPO~\cite{schulman2017proximal}.

\label{sec:theory}

\subsection{Your Language Model Is Secretly a Reward Model}

DPO circumvents both the explicit training of a reward function and the application of RL for policy improvement by reducing the process to a single maximum likelihood training step. Importantly, its objective function, given in Eq.~\ref{eq:main_eq}, is equivalent to the Bradley-Terry model where the rewards are parameterized as $r^*(x, y) = \beta \log\frac{\pi^*_\theta(y \mid x)}{\pi_\text{ref}(y \mid x)}$. The optimization thus targets the model parameters of $\pi_{\theta}$, paralleling the optimization for the reward model in Eq.~\ref{eq:reward_model} after a suitable change of variables. This section lays out the theory supporting this parameterization, demonstrates that it does not place extra restrictions on the class of reward functions that can be learned, and establishes that the optimal policy can be recovered exactly. We start by formalizing an equivalence relation among reward functions.

\begin{definition}
We define two reward functions $r(x, y)$ and $r'(x, y)$ to be equivalent if and only if ${r(x, y)-r'(x, y) = f(x)}$ for some function $f$.
\end{definition}

It is straightforward to verify that this relation is indeed an equivalence relation, resulting in a partitioning of the space of reward functions into equivalence classes. This gives rise to the following two lemmas:

\begin{lemma}\label{lemma:same_prefrence} In the context of the Plackett-Luce (and specifically the Bradley-Terry) preference model, any two reward functions within the same equivalence class will induce identical preference distributions.
\end{lemma}

\begin{lemma}\label{lemma:same_policy}
    Any two reward functions belonging to the same equivalence class will yield the same optimal policy for the constrained reinforcement learning problem.
\end{lemma}

The proofs are direct and can be found in Appendix \ref{app:lemma1}. The first lemma highlights a well-recognized under-identifiability characteristic of the Plackett-Luce class of models \cite{plackett1975analysis}. As a result, it is standard practice to introduce extra identifiability conditions to ensure meaningful maximum likelihood estimation outcomes from Eq. \ref{eq:reward_model} \cite{bong2022generalized}. The second lemma asserts that within any equivalence class, all reward functions define the same optimal policy; thus, for practical purposes, it is sufficient to recover any representative reward function from the correct equivalence class. The following theorem, proved in Appendix~\ref{app:thm1}, formalizes this fact:

\begin{theorem}\label{thm:main}
    Given mild conditions, every reward class compatible with the Plackett-Luce (and in particular the Bradley-Terry) models can be parameterized as ${r(x, y) = \beta \log \frac{\pi(y\mid x)}{\pi_\text{ref}(y\mid x)}}$, for some policy $\pi(y\mid x)$ and a chosen reference model $\pi_\text{ref}(y \mid x)$.
\end{theorem}

\begin{sproof}
    Consider an arbitrary reward function $r(x, y)$ that induces the corresponding optimal policy $\pi_r(y \mid x)$ as given by Eq. \ref{eq:op_policy}. Our aim is to demonstrate that some member of the equivalence class of $r$ admits the reparameterization claimed above. Let us introduce the projection
\begin{equation}
    f(r; \pi_\text{ref}, \beta)(x, y) = r(x, y) - \beta\log\sum_{y}\pi_\text{ref}(y\mid x)\exp\left(\frac{1}{\beta}r(x, y)\right)
\end{equation}
Here, the operator $f$ normalizes $r(x, y)$ by subtracting the log-partition function, which depends solely on $x$. As this normalization term is independent of $y$, the function $f(r; \pi_\text{ref}, \beta)(x, y)$ belongs to the same equivalence class as $r(x, y)$. Moreover, substituting $r$ with the form on the right side of Eq.~\ref{eq:main_eq}—which is valid for any reward—yields $f(r; \pi_\text{ref}, \beta)(x, y) = \beta \log \frac{\pi_r(y\mid x)}{\pi_\text{ref}(y\mid x)}$. Therefore, the operator $f$ returns an equivalent reward function with the reparameterized structure we desire. This shows that the proposed parameterization does not forfeit generality in our reward model.
\end{sproof}

Theorem~\ref{thm:main} can also be interpreted as identifying the specific reward function selected by the DPO reparameterization from each equivalence class, namely, the reward function that fulfills:

\begin{equation}\label{eq:lag_p}
     \sum_{y}\underbrace{\pi_\text{ref}(y\mid x)\exp\left(\frac{1}{\beta}r(x, y)\right)}_{=\pi(y\mid x)\text{, using Thm.~\ref{thm:main} reparam.}} = 1,
\end{equation}

meaning that $\pi(y\mid x)$ indeed forms a proper probability distribution (all values are non-negative and sum to unity). Examining Eq.~\ref{eq:op_policy}, it becomes evident that Eq.~\ref{eq:lag_p} functions as the normalization constant for the optimal policy derived from the reward $r(x, y)$. The central observation underpinning the DPO algorithm is that, by introducing constraints on the otherwise unconstrained Plackett-Luce (as well as the Bradley-Terry) class of preference models, it is possible to maintain the set of representable reward models, while simultaneously ensuring the optimal policy in Eq.~\ref{eq:op_policy} is analytically computable for any input $x$.

\subsection{Instability of Actor-Critic Algorithms} Our proposed framework further allows us to analyze and interpret the instability phenomena seen in standard actor-critic algorithms used for RLHF, such as PPO. Focusing on the RL fine-tuning phase described in Section~\ref{section:prelims}, and adopting the perspective of control as inference~\cite{levine2018reinforcement} within the context of the constrained RL problem described by Eq.~\ref{eq:RL}, we parameterize the policy as $\pi_{\theta}(y\mid x)$. The optimization objective consists of minimizing $\mathbb{D}_{\text{KL}}[\pi_{\theta}(y|x) \mid\mid \pi^*(y\mid x)]$, with $\pi^*$ denoting the optimal policy derived in Eq.~\ref{eq:optimum_model} from the reward $r_{\phi}(y, x)$. Through algebraic manipulation, we arrive at the following objective:

\begin{equation}\label{eq:AC}
    \max_{\pi_{\theta}}\mathbb{E}_{\pi_{\theta}(y\mid x)}\bigg[\underbrace{r_{\phi}(x, y) -\beta\log\sum_{y}\pi_\text{ref}(y\mid x)\exp\left(\frac{1}{\beta}r_{\phi}(x, y)\right)}_{f(r_{\phi}, \pi_\text{ref}, \beta)} - \underbrace{\beta\log\frac{\pi_{\theta}(y\mid x)}{\pi_\text{ref}(y\mid x)}}_{\text{KL}}\bigg]
\end{equation}

The objective considered here aligns with that used in previous studies 
\citep{ziegler2020finetuning, stiennon2022learning, bai2022training, ouyang2022training}, which employ a DPO-equivalent reward within the reward family defined by $r_{\phi}$. In this framework, the normalization component in $f(r_{\phi}, \pi_\text{ref}, \beta)$ can be viewed as the soft value function associated with the reference policy $\pi_\text{ref}$. Although this term leaves the set of optimal solutions unchanged, omitting it can introduce high variance in the policy gradient, thereby destabilizing training. One approach to handle the normalization is to introduce a learned value function; however, this often results in challenging optimization issues. Another common practice, found in earlier literature, is to standardize rewards using a human completion baseline—effectively a single-sample Monte Carlo approximation of the normalization. The DPO reparameterization stands in contrast, offering a reward formulation that circumvents the need for such baselines.

\section{Experiments}
Here, we provide an empirical assessment of DPO in learning policies directly from preference data. To begin, we explore a controlled text generation scenario, focusing on how effectively DPO negotiates the trade-off between maximizing reward and controlling the KL-divergence to the reference policy, in comparison to established preference learning methods like PPO. We then analyze DPO’s effectiveness on larger-scale models and more challenging RLHF benchmarks, such as summarization and dialog. Our results indicate that, with minimal hyperparameter tuning, DPO frequently matches or outperforms robust baselines such as PPO-based RLHF and outstrips strategies that select the best out of $N$ sampled rollouts using a learned reward model. Before detailing these results, we outline the experimental protocol; further specifics can be found in Appendix~\ref{app:exp_details}.

\textbf{Tasks.} The suite of experiments encompasses three distinct open-ended text generation tasks. For each task, learning algorithms are tasked with optimizing a policy based on a dataset of preference comparisons, denoted as $\mathcal{D}=\bigl\{x^{(i)}, y_w^{(i)}, y_l^{(i)}\bigr\}_{i=1}^N$. In the context of \textbf{controlled sentiment generation}, the input $x$ represents a prompt extracted from a movie review within the IMDb dataset \cite{maas-EtAl:2011:ACL-HLT2011}. Here, the policy is responsible for generating a completion $y$ characterized by positive sentiment. To systematically assess performance, preference pairs over generated texts are produced using an off-the-shelf sentiment classifier, ensuring that $p(\text{positive}\mid x,y_w)>p(\text{positive}\mid x,y_l)$. For SFT, we further adapt GPT-2-large through fine-tuning on the training partition of IMDb reviews, continuing until model performance stabilizes (additional details in App~\ref{app:sentiment_details}). For the \textbf{summarization} task, $x$ is drawn from Reddit forum posts, and the objective is to produce a summary $y$ capturing the main ideas. Building upon established benchmarks, we employ the Reddit TL;DR summarization dataset \citep{volske-etal-2017-tl}, integrating human-annotated preference data from \citeauthor{stiennon2022learning}. We start from an SFT model specifically adapted to human-crafted summaries of forum posts\footnote{\url{https://huggingface.co/CarperAI/openai_summarize_tldr_sft}}, leveraging the TRLX framework \citep{leandro_von_werra_2023_7790115} for RLHF. The preference annotations are sourced from the work of \citeauthor{stiennon2022learning}, applied to completions generated by an alternative SFT model trained in a comparable fashion. Lastly, for \textbf{single-turn dialogue}, the prompt $x$ consists of a user's question, which can range from scientific inquiries to requests for personal guidance. Here, the policy must craft an informative and engaging reply $y$. This task uses the Anthropic Helpful and Harmless dialogue dataset \citep{bai2022training}, comprised of 170k dialogues between a human interlocutor and a virtual assistant. Each dialogue concludes with two responses from a large, unspecified language model, one of which is labeled as preferred by a human annotator. Since there is no publicly available SFT model for this dataset, we generate one by fine-tuning a readily available language model exclusively on the responses favored in human evaluations.

\textbf{Evaluation.} We adopt two distinct strategies for evaluating our methods. For assessing each algorithm’s capability to optimize the constrained reward maximization objective within the controlled sentiment generation setting, we plot the Pareto frontier illustrating the trade-off between obtained reward and KL-divergence from the reference policy. Because the underlying reward function (sentiment classifier) is accessible, the exact frontier is calculable in this scenario. Conversely, in practical deployments where the true reward function is unknown, algorithmic performance is measured by their \textit{win rate} relative to a baseline policy. Here, GPT-4 acts as a surrogate for human judgment, evaluating summary quality and response helpfulness in the summarization and single-turn dialogue scenarios, respectively. The summarization baseline consists of reference summaries from the test set; in dialogue, the preferred test set responses are used as the baseline. While prior literature has found that LMs often outperform standard metrics as automated evaluators \citep{Chen2023ExploringTU}, we corroborate our reliance on GPT-4 in Sec.~\ref{sec:human-judgments} by conducting an additional human subject study. Our results reveal a strong alignment between GPT-4 assessments and those of human annotators, with levels of agreement matching or even exceeding inter-annotator consensus.

\textbf{Methods.} Alongside DPO, we benchmark several existing techniques designed to align language models with human preferences. The most straightforward baselines involve zero-shot prompting with \textbf{GPT-J} \citep{gpt-j} for summarization and two-shot prompting using \textbf{Pythia-2.8B} \citep{biderman2023pythia} for dialogue. We also evaluate the \textbf{SFT} model, as well as \textbf{Preferred-FT}—a variant finetuned through supervised learning to prefer the selected output $y_w$ generated either by the SFT model (in both controlled sentiment and summarization experiments) or a generic language model (in single-turn dialogue). Another related method is \textbf{Unlikelihood}~\citep{welleck2019neural}, which directly maximizes the policy's likelihood for $y_w$ and explicitly reduces the likelihood for $y_l$, employing a tunable coefficient $\alpha\in[0,1]$ for the “unlikelihood” component. Our suite of baselines also includes \textbf{PPO} \citep{schulman2017proximal}, which leverages a reward model trained on preference data, and \textbf{PPO-GT}, an oracle variant that can optimize using the true reward function in controlled sentiment scenarios. For PPO-GT, we run both a standard implementation \cite{leandro_von_werra_2023_7790115} and a custom version featuring reward normalization and hyperparameter refinement for enhanced effectiveness (these improvements are likewise used for ‘standard’ PPO with learned rewards). Finally, we test the \textbf{Best of $N$} baseline, which samples $N$ candidate completions from the SFT model (or Preferred-FT for dialogue) and selects the highest-rewarding response according to a preference-trained reward model. While this approach can achieve high performance independent of PPO’s reward model integration, it is often impractical for nontrivial $N$ due to the need to generate $N$ responses per test query.

\subsection{How well can DPO optimize the RLHF objective?}

In standard RLHF algorithms, the objective of maximizing reward while imposing a KL constraint serves to find a balance between exploiting high-reward behaviors and ensuring that the policy does not stray excessively from the reference policy. As such, evaluating and comparing algorithms requires considering both the achieved reward and the magnitude of the KL divergence; gaining marginally more reward at the expense of a significantly increased KL is often suboptimal. In Figure~\ref{fig:frontier-tldr-main}, we present the reward versus KL tradeoff across different methods in the sentiment domain. For each algorithm, we conduct several training runs, varying the policy conservativeness parameter in each case (specifically, target KL $\in\{3,6,9,12\}$ for PPO, $\beta \in \{0.05,0.1,1,5\}$, and $\alpha\in\{0.05,0.1,0.5,1\}$ for unlikelihood, with random seeds employed for preferred-FT), leading to a total of 22 experiment runs. At every interval of 100 training steps, and continuing until the algorithm reaches convergence, we assess each resulting policy on a test prompt set, recording both the mean reward according to the ground-truth reward function and the average sequence-level KL\footnote{Defined here as the aggregate sum of KL divergences over all timesteps in the sequence.} against the reference policy, $\text{KL}\left(\pi\mid \mid \pi_\text{ref}\right)$. Our findings indicate that DPO consistently delivers the most favorable reward-KL tradeoff, achieving superior rewards with lower KL values compared to alternatives. This outcome stands out for several reasons. Notably, DPO and PPO seek to optimize an identical objective; yet, DPO substantially outperforms PPO in terms of reward efficiency at any given KL value, fully dominating PPO on this frontier. Additionally, DPO continues to surpass PPO's performance \emph{even in scenarios where PPO is provided access to true reward signals} (PPO-GT).

\subsection{Can DPO scale to real preference datasets?}
\label{sec:dpo-real-datasets}

We proceed to assess the effectiveness of DPO fine-tuning in the context of summarization and single-turn dialogue tasks. Notably, established automated metrics for summarization like ROUGE are known to show weak alignment with human judgments~\citep{stiennon2022learning}. Previous studies demonstrate that adapting language models with PPO using human feedback results in summaries that better reflect human preferences. In our evaluation, we generate outputs on the TL;DR summarization dataset's test split, then measure the mean win rate relative to the reference completions present in the test set. Output samples for every approach are generated over a temperature range from 0.0 up to 1.0, and the resulting win rates are depicted in Figure~\ref{fig:frontier-tldr-main} (right). DPO, PPO, and Preferred-FT each use the identical GPT-J SFT model for fine-tuning\footnote{\url{https://huggingface.co/CarperAI/openai_summarize_tldr_sft}}. Our results indicate that DPO attains an average win rate close to 61\% at temperature 0.0, outperforming PPO, which achieves approximately 57\% at its optimal (0.0) temperature. Additionally, DPO surpasses the best $N$ baseline in terms of maximum win rate achieved. It is important to mention that we did not systematically optimize the DPO $\beta$ hyperparameter; thus, these figures may not fully capture the algorithm's capabilities. Furthermore, DPO demonstrates a greater resilience to temperature changes, whereas PPO's performance can regress toward the original GPT-J model's results at higher temperatures. Preferred-FT yields minimal gains compared to the SFT baseline. A direct human preference comparison between DPO and PPO, detailed in Section~\ref{sec:human-judgments}, reveals that DPO samples at temperature 0.25 are selected 58\% of the time over PPO samples at the same temperature.

For single-turn dialogue, we assess various approaches using the portion of the Anthropic HH dataset \citep{bai2022training} test set that includes a single exchange between a human and assistant. To determine win rates, GPT-4 comparisons use the preferred test completions as references across different techniques. Since there is no standard SFT model for this setting, our process begins with the pre-trained Pythia-2.8B model; we apply Preferred-FT on the selected completions to fit a reference model, ensuring the generated completions match the training distribution, and proceed to DPO training. Additionally, we compare performance to the best among 128 Preferred-FT completions (this plateaued at $N=128$, see Appendix Figure~\ref{fig:best-of-n}) as well as to a 2-shot prompt-based variant of Pythia-2.8B. Across the best-performing sampling temperatures, DPO matches or exceeds competing methods. We also test an RLHF model using PPO, fine-tuned on the Anthropic HH dataset \footnote{\url{https://huggingface.co/reciprocate/ppo_hh_pythia-6B}} and sourced from a reputable repository \footnote{\url{https://github.com/CarperAI/trlx/tree/main/examples/hh}}, but are unable to identify a prompting strategy or temperature that surpasses base Pythia-2.8B performance. Our prior findings from TL;DR, along with the fact that both approaches aim to maximize the same reward function, lead us to treat Best of 128 as a rough indicator of PPO capability. In summary, DPO stands out as the only approach that is computationally efficient while also consistently outperforming preferred completions in the Anthropic HH dataset, matching or bettering the results of the more resource-intensive Best of 128 baseline. Lastly, Figure~\ref{fig:dialogue-main} illustrates that DPO reaches optimal performance after comparatively few training steps.

\subsection{Generalization to a new input distribution}

To further assess the robustness of PPO and DPO in the face of distributional changes, we evaluate the trained policies from our Reddit TL;DR summarization study on an out-of-distribution dataset—specifically, on news articles from the test partition of the CNN/DailyMail dataset \citep{nallapati-etal-2016-abstractive}. We apply the best sampling temperatures identified for TL;DR (0 and 0.25) to both policies and report the outcomes in Table~\ref{tab:ood}. Using the GPT-4 win rate as an evaluation metric, and leveraging the same GPT-4 (C) prompt from the Reddit TL;DR study (with the phrase ``forum post'' replaced by ``news article''), we compare generated summaries to ground-truth references. On this new data distribution, DPO maintains a substantial advantage over PPO. These findings offer preliminary support that DPO policies are capable of generalizing to novel input distributions on par with or exceeding PPO, despite DPO not incorporating additional unlabeled Reddit TL;DR prompts as PPO does.

\subsection{Validating GPT-4 judgments with human judgments}
\label{sec:human-judgments}
We design a human evaluation to corroborate the validity of GPT-4’s assessments, focusing on outcomes from the TL;DR summarization experiment and employing two distinct GPT-4 prompt formats. The \textbf{GPT-4 (S)} (simple) prompt asks the model to choose the summary that better captures the main points of the post. The \textbf{GPT-4 (C)} (concise) prompt also requests consideration of conciseness, a factor included because GPT-4, when given the simple prompt, tends to favor longer, more redundant summaries than do human judges. Complete prompt wordings can be found in Appendix~\ref{app:prompts}. In the study, we contrast three methods, spanning a range of expected summary quality: the top-performing DPO (temp. 0.25), the lowest-performing PPO (temp. 1.0), and an intermediate SFT (temp. 0.25), each evaluated against greedily-sampled PPO (its optimal temperature). With both prompt types, we observe that GPT-4’s preferences align with human judgment at rates similar to human-human agreement, implying that GPT-4 is an effective stand-in for human evaluation (note: multiple human ratings are collected only for DPO and PPO-1 comparisons due to the number of available raters). Generally, the \textbf{GPT-4 (C)} prompt yields win rates most consistent with human evaluations, motivating its use in Section~\ref{sec:dpo-real-datasets} as the primary assessment tool. Further information regarding the human evaluation protocol, including the web interface for raters and the pool of volunteer participants, is detailed in Appendix~\ref{app:human-study}.

\section{Discussion}
Preference-based learning represents a robust and extensible approach for developing language models that are both effective and aligned with human values. In this work, we proposed DPO—a straightforward training procedure designed to leverage human preference data for language model training, eschewing reinforcement learning entirely. Instead of forcing preference learning into the conventional RL framework to utilize pre-existing RL methods, DPO establishes an explicit correspondence between language model policies and reward functions. This enables language models to be optimized with respect to human preferences \textit{directly}, employing a standard cross-entropy objective, and avoids both the use of reinforcement learning and any compromise in expressiveness. DPO demonstrates performance on par with, or exceeding, leading RLHF techniques such as those relying on PPO, and achieves this with almost negligible hyperparameter adjustment. As a result, DPO provides a significant step toward lowering the entry barriers for preference-driven language model training.

\textbf{Limitations \& Future Work.} Our findings prompt several avenues for future investigation. Notably, questions remain regarding the generalization capacity of DPO-trained policies, especially in comparison to models trained on explicit reward signals when confronted with out-of-distribution data. Preliminary experiments indicate that the generalization abilities of DPO policies are at least comparable to those derived from PPO, though further systematic exploration is warranted. Additionally, it would be worthwhile to determine whether employing DPO-based self-labeling techniques could facilitate more effective use of prompts lacking manual annotation. Another pertinent line of inquiry pertains to the manifestation of reward over-optimization within the direct preference optimization context, and whether the modest performance decline observed in Figure~\ref{fig:dialogue-main}-right exemplifies this phenomenon. While our experiments utilize models up to 6B parameters, investigating the scalability of DPO to significantly larger, state-of-the-art architectures is a compelling future direction. In terms of assessment methodology, we observed that GPT-4's computed win rates can be sensitive to prompt formulation; determining optimal strategies for eliciting reliable automated evaluations remains an open question. Lastly, DPO holds promise for a wide range of applications outside preference-based language model training, extending potentially to the training of generative models across diverse data modalities.

\end{document}